% CODEX REPORT::LEAD::FINAL-LINK-ARC::2026-10-08
\documentclass[11pt]{article}
\usepackage[T1]{fontenc}
\usepackage[margin=0.82in]{geometry}
\usepackage{amsmath,amssymb,amsthm,booktabs,array,mathtools,xcolor,xurl,hyperref}
\hypersetup{colorlinks=true,linkcolor=blue!55!black,citecolor=blue!55!black,urlcolor=blue!55!black}
\setlength{\emergencystretch}{3em}
\setlength{\parskip}{0.35em}
\setlength{\parindent}{0pt}
\newtheorem{theorem}{Theorem}[section]
\newtheorem{lemma}[theorem]{Lemma}
\newtheorem{corollary}[theorem]{Corollary}
\newtheorem{definition}[theorem]{Definition}
\newtheorem{proposition}[theorem]{Proposition}
\newcommand{\R}{\mathcal R}
\newcommand{\B}{\mathcal B}
\newcommand{\SAT}{\operatorname{SAT}}
\newcommand{\poly}{\operatorname{poly}}
\newcommand{\eps}{\varepsilon}
\newcommand{\st}{\textnormal{s.t.}}
\title{FINAL-LINK-ARC:\\The Weakest Charged Successor Needed for Finite Entry into the Williams Region}
\author{Lead ARC audit}
\date{8 October 2026}

\begin{document}
\maketitle

\begin{abstract}
The fixed wrapper cell with fanout exponent $1/3$ and child-size exponent
$2/3$ is sufficient, but it is not necessary.  The weakest proof-spine
statement currently isolated is a cofinal finite-hit theorem: for a cofinal
set of possible polynomial SAT exponents, one needs one finite,
source-checkable route from the corresponding achievable seed to any
model-matched forbidden point.  A reusable sufficient form is a charged
max-affine successor.  If a compiler maps a SAT algorithm with exponent pair
$(t,s)$ to one with
\[
 t'\leq \max\{r,\max_j(w_j+\lambda_jt)\},\qquad
 s'\leq \max\{p,\max_j\lambda_js\},
\]
where every $\lambda_j<1$, then the limiting envelope is
\[
 C_* = \max\{r,\max_j w_j/(1-\lambda_j)\},\qquad S_*=p.
\]
The final link closes whenever $C_*+p<2\cos(\pi/7)$ and the inequalities are
witnessed by one uniform exact machine at every finite return.  Existing
recovered wrapper theorems verify parts of this accounting, but none proves a
strict all-input general-SAT successor.  In particular, PRF3-0079 invalidates
the earlier half-scale pointwise acquisition cell for an all-block schedule:
the omitted fanout changes the exponent from $\beta+\lambda k$ to
$1-\lambda+\beta+\lambda k$.  Thus P versus NP remains \textsc{Open}.
\end{abstract}

\paragraph{Status.}
The Williams endpoint and the finite-entry deductions in this report are
\textsc{Known Theorem} and \textsc{Proved/Exact}, respectively.  The numerical
audit is \textsc{Computational}.  The existence of a general-SAT compiler or
finite portfolio satisfying the charged inequalities is \textsc{Open} and is
the only load-bearing constructional arrow discussed here.  No P versus NP
proof is claimed.

\section{Endpoint and model}

Let $N$ denote the encoded input length of a CNF instance.  Define the
achievable region $\R$ to contain $(t,s)$ when one deterministic, uniform,
exact random-access SAT algorithm has simultaneous bounds
\[
 T(N)\leq N^{t+o(1)},\qquad S(N)\leq N^{s+o(1)}.
\]
The same machine must witness both coordinates.  The region is upward closed.
Set
\[
 \gamma:=2\cos(\pi/7)=1.8019377358\ldots,
 \qquad
 \B:=\{(t,s):t+s<\gamma\}.
\]

Williams proves the following product lower bound in the deterministic
random-access model \cite[Theorem 3.9]{Williams2010}.

\begin{theorem}[Williams product endpoint]
If $t(N)$ and $s(N)$ are polynomially bounded and an algorithm solves SAT in
time $t(N)$ and space $s(N)$, then for every $\eps>0$,
\[
 t(N)s(N)=\Omega\!\left(N^{\gamma-\eps}\right).
\]
Consequently, no pair in $\R$ lies in $\B$.
\end{theorem}

Indeed, a machine with exponent pair $(t,s)$ and $t+s<\gamma$ contradicts the
theorem after choosing $\eps$ smaller than the positive gap.  This strict
margin also absorbs the finitely many $o(1)$ terms introduced below.  A
different SAT representation may be used only after charging its exact length
change and its simulation overhead in the same machine model.

\section{The logically weakest finite-entry statement}

The final link does not logically require one universal contraction factor, a
fixed route length, or the particular point $(1,2/3)$.  The following is the
cofinal formulation already implicit in the strongest current proof spine.

\begin{definition}[Source-checkable finite route]
A finite route is source-checkable when every vertex is witnessed by one
uniform exact SAT machine, every edge is justified by a finite machine
construction or by an explicit existence theorem, and the proof charges
construction, output, recursive calls, replay, retained state, and workspace.
The route depth is a constant after the starting machine is selected; it may
not grow with $N$.
\end{definition}

\begin{theorem}[Cofinal charged hit]
Let $\mathcal K\subseteq\mathbb N$ be unbounded.  Suppose that for every
$K\in\mathcal K$, whenever $(K,K)\in\R$, there is a source-checkable finite
route in $\R$ from $(K,K)$ to $\B$.  Then $P=NP$ contradicts the Williams
endpoint.
\end{theorem}

\begin{proof}
Under $P=NP$, some deterministic exact SAT machine runs in time $N^{K_0}$ for
a finite $K_0$.  Its workspace is at most its time after harmless adjustment,
so $(K_0,K_0)\in\R$.  Choose $K\in\mathcal K$ with $K\geq K_0$.  Upward
closure gives $(K,K)\in\R$, and the assumed finite route reaches $\B$, which
the endpoint forbids.
\end{proof}

This theorem is weaker than a compiler applying to every achievable pair.  It
allows the successful route, its contraction factors, and its finite length
to depend on $K$.  Cofinality replaces advance knowledge of the unknown
polynomial exponent.  It does not permit input-length-dependent advice,
input-by-input algorithm choice, or a growing recursion depth.

\section{A reusable max-affine successor theorem}

The most useful source-checkable target is stronger than cofinal charged hit,
but substantially weaker than the original fixed cell.

\begin{theorem}[Charged max-affine successor]
Assume there is a fixed finite compiler which transforms every uniform exact
SAT machine with achievable pair $(t,s)$ into a uniform exact SAT machine with
pair $(t',s')$ satisfying
\begin{align}
 t'&\leq f(t):=\max\{r,\max_{j\in J}(w_j+\lambda_jt)\},\label{eq:maxafftime}\\
 s'&\leq g(s):=\max\{p,\max_{j\in J}\lambda_js\},\label{eq:maxaffspace}
\end{align}
where $J$ is finite, $r,p,w_j\geq0$, and $0\leq\lambda_j<1$.  All constants
are independent of the input and the supplied machine.  Define
\[
 \lambda_{\max}:=\max_j\lambda_j,qquad
 C_*:=\max\left\{r,\max_j\frac{w_j}{1-\lambda_j}\right\}.
\]
Then repeated finite returns satisfy
\begin{align*}
 t_m&\leq C_*+\lambda_{\max}^m(t_0-C_*)_+,\\
 s_m&\leq p+\lambda_{\max}^m(s_0-p)_+.
\end{align*}
If
\begin{equation}
 \boxed{C_*+p<\gamma,}\label{eq:maincriterion}
\end{equation}
then every finite starting pair enters $\B$ after finitely many returns.
Consequently, the compiler and the Williams endpoint imply $P\ne NP$.
\end{theorem}

\begin{proof}
For every channel $j$, the definition of $C_*$ gives
$w_j+\lambda_jC_*\leq C_*$, and $r\leq C_*$.  Thus $f(C_*)\leq C_*$.  On the
other hand, every $x$ with $f(x)\leq x$ must satisfy $x\geq r$ and
$x\geq w_j/(1-\lambda_j)$ for every $j$, so $C_*$ is the least pre-fixed
point.  The map $f$ is $\lambda_{\max}$-Lipschitz.  Monotonicity and the
pre-fixed-point property imply
\[
 f^m(t_0)\leq C_*+\lambda_{\max}^m(t_0-C_*)_+.
\]
The identical argument for $g$ gives limiting envelope $p$.  If
$\Delta=\gamma-(C_*+p)>0$, choose finite $m$ so that the two transient terms
sum to less than $\Delta$.  The $m$-fold compiled machine has one simultaneous
pair in $\B$.  Under $P=NP$ a finite starting pair exists, so the endpoint is
contradicted.
\end{proof}

The compiler need not discover the unknown starting exponents.  The
meta-proof selects a finite return count after selecting the hypothetical
polynomial SAT machine.  The source of the final machine is a finite
composition and therefore remains uniform.  If an operation in the compiler
calls the supplied solver, its cost belongs in a term $w+\lambda t$; it may not
be hidden inside the constant $r$.  A parent-scale call has $\lambda=1$ and
does not satisfy the theorem.

\section{Complete one-channel ledger}

Suppose one compiler invocation has the following charged parameters:
\begin{description}
\item[$a$] at most $N^{a+o(1)}$ recursively relevant children;
\item[$b$] each recursive child has encoded or logical source size at most
$N^{b+o(1)}$, with $0\leq b<1$;
\item[$\beta$] an additional per-child interface or table multiplicity
$N^{\beta+o(1)}$;
\item[$g$] per-call generation, virtual-access, or replay overhead
$N^{g+o(1)}$;
\item[$r$] the maximum exponent of all nonrecursive discovery, construction,
materialization, exact verification, combination, and shared setup;
\item[$p$] peak live workspace for the construction, interfaces, handles,
generators, retained source state, and ordinary child buffers.
\end{description}
If every recursive call uses the supplied machine on a size-$N^b$ child, the
successor recurrence is
\begin{align}
 t'&\leq\max\{r,a+\beta+g+bt\},\label{eq:onechanneltime}\\
 s'&\leq\max\{p,bs\}.\label{eq:onechannelspace}
\end{align}
Hence
\begin{equation}
 \boxed{
 C_* = \max\left\{r,\frac{a+\beta+g}{1-b}\right\},\quad
 S_*=p,\quad
 \max\left\{r,\frac{a+\beta+g}{1-b}\right\}+p<\gamma.}
 \label{eq:onechannelcriterion}
\end{equation}

For ordinary explicit children, materializing all output forces the charged
time guard $r\geq a+b$ in the direct schedule, and holding one child normally
forces $p\geq b$.  With local workspace exponent $s_0$, one may write
$p=\max\{s_0,b\}$.  Setting $\beta=g=0$ recovers the current one-type wrapper
formula
\[
 C_* = \max\{r,a/(1-b)\},\qquad
 S_* = \max\{s_0,b\}.
\]

If $N^a$ explicit children of size $N^b$ cover $\Theta(N)$ distinct source
items, elementary incidence counting gives $a+b\geq1$.  At the boundary
$a=1-b$, with linear construction and no further overhead, $C_*=1$.  The
Williams condition is then merely
\[
 b<\gamma-1=0.8019377358\ldots.
\]
Thus $(a,b)=(1/3,2/3)$ is one convenient witness, not a unique target.

\begin{center}
\begin{tabular}{lccccc}
\toprule
explicit cell $(a,b)$ & $r$ & $s_0$ & $C_*$ & $S_*$ & $\gamma-(C_*+S_*)$\\
\midrule
$(1/3,2/3)$ & $1$ & $2/3$ & $1$ & $2/3$ & $0.135271\ldots$\\
$(1/2,1/2)$ & $1$ & $1/2$ & $1$ & $1/2$ & $0.301938\ldots$\\
$(2/5,3/5)$ & $1$ & $1/5$ & $1$ & $3/5$ & $0.201938\ldots$\\
$(1/4,3/4)$ & $1$ & $1/2$ & $1$ & $3/4$ & $0.051938\ldots$\\
\bottomrule
\end{tabular}
\end{center}

More generally, the explicit one-type region is exactly
\[
 \max\{r,a/(1-b)\}+\max\{s_0,b\}<\gamma.
\]
Equivalently, after fixing $b$ and $s_0$, it is enough that
\[
 r<\gamma-\max\{s_0,b\},\qquad
 a<(\gamma-\max\{s_0,b\})(1-b),
\]
with both strict inequalities interpreted together through the maximum.

\section{PRF3-0079 and the corrected Route B arithmetic}

The recovered Route B cell used a selected $N^{k+o(1)}$ SAT witness to fill
$N^{\beta+o(1)}$ interface entries for one source-localized child of size
$N^{\lambda+o(1)}$.  Its one-child acquisition exponent is correctly
\[
 e_{\mathrm{one}}=\beta+\lambda k.
\]
Combined with the other D1 table constraint, this motivated the concrete
laboratory point $k=2$, $C=4/3$, $\lambda=1/2$, $\beta<1/6$.

PRF3-0079 proves that this is not the cost of an all-block descendant
algorithm.  If exact coverage requires $N^{a+o(1)}$ blocks and there is no
cross-block sharing, then
\begin{equation}
 e_{\mathrm{global}}=a+\beta+\lambda k.\label{eq:globalacq}
\end{equation}
For the natural explicit all-source cover, $a=1-\lambda$.  At $k=2$ and
$C=4/3$, the target $e_{\mathrm{global}}<C$ becomes
\[
 1-\lambda+\beta+2\lambda<\frac43
 \quad\Longleftrightarrow\quad
 \boxed{\lambda<\frac13-\beta.}
\]
Consequently, $\beta$ near $1/6$ requires $\lambda<1/6$, while
$\beta\to0$ permits only $\lambda<1/3$.  The former half-scale target is
\textsc{Falsified/Retired} for this all-block pointwise schedule.  It remains
available only to a different CoverageContract which proves that one selected
child suffices, or to a proved shared BatchAcquire schedule.

The same correction can be tested directly against the Williams product
endpoint.  Allowing virtual-access overhead $g$, the all-block successor is
\begin{align}
 t'&\leq\max\{r,1-\lambda+\beta+g+\lambda t\},\label{eq:routeBtime}\\
 s'&\leq\max\{p,\lambda s\}.
\end{align}
Its envelope is
\begin{equation}
 C_* = \max\left\{r,1+\frac{\beta+g}{1-\lambda}\right\},
 \qquad S_*=p.\label{eq:routeBfixed}
\end{equation}
For half-scale ordinary explicit children, $p\geq1/2$.  If $r$ does not
dominate and $g=0$, Williams requires
\[
 1+2\beta+\frac12<\gamma
 \quad\Longleftrightarrow\quad
 \boxed{\beta<\frac{\gamma-3/2}{2}=0.1509688679\ldots.}
\]
Thus $\beta<1/6$ is not sufficient for the explicit half-scale implementation.
The rational choice $\beta=1/7$ gives
\[
 (C_*,S_*)=(9/7,1/2),\qquad C_*+S_*=25/14<\gamma.
\]
This is a valid conditional numerical target, not a construction.

\subsection{Ordinary children versus source-queryable virtual children}

An ordinary child is a finite SAT input that is explicitly written and then
read by the recursive solver.  Its output traffic belongs in $r$, and its live
buffer belongs in $p$.  A short label naming such a child does not remove
either charge.

A virtual child is different.  It consists of a handle and a proved random
access decoder that returns every requested bit of the logical child from the
parent source.  In that model, $p$ need not include $\lambda$, but all decoder
time must appear in $g$, and $p$ must include the handle, decoder state,
retained source state, recursion stack, and any cache needed for exact access.
The recursive machine model must permit the simulation without silently
granting a second read-only input or free oracle access.

Conditionally, if $\lambda=1/2$, $\beta=1/6$, $g=0$ in exponent, and
$p\leq1/6$, then the virtual envelope is
\[
 (C_*,S_*)=(4/3,1/6),\qquad C_*+S_*=3/2<\gamma.
\]
This explains why the older half-scale numbers can reappear for a genuine
source-queryable representation.  What remains open is precisely the decoder
and coverage theorem that makes those numbers honest on every input.

\subsection{What a global shared object \texorpdfstring{$H$}{H} must accomplish}

Suppose $H$ is constructed once in time exponent $h$ and replaces the
$N^\beta$ pointwise acquisition calls.  If recursive children still number
$N^a$ and each call has access/recovery overhead $N^g$, the recurrence becomes
\[
 t'\leq\max\{h,r_0,a+g+\lambda t\}.
\]
Therefore
\begin{equation}
 C_* = \max\left\{h,r_0,\frac{a+g}{1-\lambda}\right\}.\label{eq:batchfixed}
\end{equation}
For $a=1-\lambda$ and $g=0$, this is $\max\{h,r_0,1\}$.  At half scale with
explicit children, Williams requires
\[
 \max\{h,r_0,1\}+\frac12<\gamma,
\]
so in particular $h,r_0<1.3019377358\ldots$.  A bound merely below $4/3$ is
not sufficient for this explicit-space implementation.  With a verified
virtual interface of space at most $1/6$, the weaker design target $h<4/3$
would be sufficient, because $4/3+1/6=3/2<\gamma$.

The object $H$ must jointly provide exact coverage, all boundary answers,
child/source representation, access, and recursive composability.  A global
source plus short handles is not enough when aggregate source traffic or
per-access replay restores the omitted exponent.

\section{Finite portfolios and disjunctive wrappers}

A finite portfolio can be valuable because different exact mechanisms may
cover different structural cases.  It does not remove worst-case charges.
Suppose portfolio member $j$ either returns a recognizable failure or an exact
SAT-preserving decomposition with channel $(w_j,\lambda_j)$.  Assume the
portfolio is finite, every input is covered by at least one member, and all
failed attempts have charged time and space included in $r$ and $p$.  Running
the portfolio sequentially yields exactly the max-affine successor
\[
 t'\leq\max\{r,\max_j(w_j+\lambda_jt)\},\qquad
 s'\leq\max\{p,\max_j\lambda_js\}.
\]
Thus its criterion is still Equation~\eqref{eq:maincriterion}.  The maximum is
coordinatewise because worst-case time and worst-case space may occur on
different inputs.  Average success or good behavior on sampled formulas does
not establish the needed worst-case successor.

A disjunctive wrapper is exact when it proves
\[
 \SAT(F)=\bigvee_{i=1}^{q(F)}\SAT(G_i).
\]
It is then simply a recursive fanout channel, and $q(F)$ contributes to $a$.
Running parallel wrappers $W_A(F),W_B(F),\ldots$ helps only if their combined
failure/success contract covers every input and their total fanout,
construction, and workspace satisfy the same ledger.  Choosing the successful
wrapper using a parent-scale SAT answer makes the selector circular and
introduces a $\lambda=1$ term.

\section{Typed return cycles}

One global contraction on every edge is stronger than necessary.  Let $I$ be
a finite set of algorithm/interface types.  For an edge $e:i\to j$, let
$\theta_e\in[0,1]$ be the source scale, let
$w_e=\eta_e+\beta_e+g_e$ charge fanout and access, and let $h_e$ charge peak
edge state.  Let $r_i$ and $p_i$ be local time and space charges.  The exact
Bellman ledger is
\begin{align}
 C_i&\geq r_i,& C_i&\geq w_e+\theta_eC_j &&(e:i\to j),\label{eq:belltime}\\
 \rho_i&\geq p_i,& \rho_i&\geq h_e,&
 \rho_i&\geq\theta_e\rho_j &&(e:i\to j).\label{eq:bellspace}
\end{align}

\begin{theorem}[Finite typed return-cycle criterion]
Assume every reachable directed cycle $Z$ satisfies
\[
 \Theta_Z:=\prod_{e\in Z}\theta_e<1.
\]
Assume also exact semantics, finite branching, and source-checkable charges on
every edge.  Then the least finite budgets solving
Equations~\eqref{eq:belltime} and \eqref{eq:bellspace} exist.  The time budget
has the path representation
\[
 C_i^*=\sup_{P=e_0\cdots e_{m-1}}
 \left[
 \sum_{k=0}^{m-1}\left(\prod_{\ell<k}\theta_{e_\ell}\right)w_{e_k}
 +\left(\prod_{\ell<m}\theta_{e_\ell}\right)r_{\operatorname{end}(P)}
 \right].
\]
If the root type $0$ satisfies $C_0^*+\rho_0^*<\gamma$, then finite unrolling
from every finite starting exponent reaches the Williams forbidden region.
\end{theorem}

\begin{proof}
Because the graph is finite and every cycle contracts, products along paths
of increasing length tend uniformly to zero.  Repeated cycles contribute
convergent geometric sums to the displayed path value, so the supremum is
finite and is the least solution of the max-affine Bellman inequalities.  The
space equations have the same monotone path interpretation.  An unknown leaf
exponent is multiplied by the path product and therefore becomes smaller than
any fixed positive margin at a finite depth.  Choose that finite depth after
the starting machine is fixed.  The resulting finite inlining is one uniform
machine, and the strict root inequality invokes the Williams endpoint.
\end{proof}

A same-scale edge $\theta_e=1$ is therefore allowed when every complete return
cycle containing it has product less than one.  A reachable unit-product cycle
does not suppress the unknown leaf exponent and cannot eliminate the unknown
$K$ by finite unrolling.  This typed theorem is the correct home for a double
wrapper whose outer phase has no contraction but whose full inner-outer cycle
does.

\section{What the recovered exact components do and do not prove}

The repository and recovered Drive audit contain valuable exact components.
Their strongest safe composition is as follows.

\paragraph{PRF3-0016 and PRF3-0018: \textsc{Proved/Exact}.}
They give recurrence arithmetic and rate-controlled uniformization after a
shrinker or ExtractorSpec is supplied.  They do not construct the shrinker.

\paragraph{PRF3-0079: \textsc{Proved/Exact}.}
It restores the all-block fanout term in pointwise acquisition.  It corrects,
rather than closes, Route B.

\paragraph{Recovered recurrence ledgers: \textsc{Proved/Exact} accounting.}
MULTICHANNEL1, STREAMCOMP1, CYCLEGAIN1, and ORBITBATCH1 support
Equations~\eqref{eq:onechannelcriterion},
\eqref{eq:batchfixed}, and the typed Bellman ledger.  Their hypotheses still
require an exact low-cost generator, source interface, or cover.

\paragraph{WRAPPED-SELFRED1 and recursive path exposure: \textsc{Proved/Exact}
semantics.}
The exposed paths use the current-scale SAT answer or do not prove a favorable
all-input resource successor.  Exactness alone is not contraction.

\paragraph{Canonicalization and decision-then-emit: \textsc{Falsified/Retired}
as independent leverage.}
Under $P=NP$ they can emit a tiny canonical answer instance, but obtaining it
already pays the unknown $N^K$ decision cost.  The resulting successor is an
identity or worse map in the charged ledger.

\paragraph{Bounded-width frontier dynamic programming: \textsc{Proved/Exact}
on a restricted class.}
Branch-width and treewidth methods provide genuine exact interfaces when the
width is small \cite{AlekhnovichRazborov2002,AllenderEtAl2014}.  No theorem in
the recovered corpus produces such width, or a comparably small backdoor, for
every CNF at favorable total cost.

\paragraph{Hashes, prime labels, random bitstreams, and route fingerprints:
\textsc{Proved/Exact} identity or replay aids.}
They distinguish histories and can protect exact cache reuse.  Changing a
four-digit hash to six or eight digits changes collision and storage costs,
not the number or encoded size of exact semantic classes.

\paragraph{Double wrappers, portfolios, and disjunctive variants:
\textsc{Conditional}.}
Their semantic composition can be exact, but the current corpus supplies no
all-input cover with the strict root budget required above.

It follows that the recovered exact components currently imply no strict
all-input successor point for general SAT in $\R$.  They imply restricted
successors, exact bookkeeping conditionals, and correctness infrastructure.
None supplies the missing general-SAT generator or cover.  This is the key
distinction between having addressed many subproblems and having closed the
final link.

\section{Literature transfer audit}

Williams's product theorem is a directly transferable endpoint in the stated
random-access model \cite{Williams2010}.  Fortnow, Lipton, van Melkebeek, and
Viglas provide an earlier family of deterministic random-access SAT
time-space exclusions below the golden-ratio time exponent
\cite{FortnowEtAl2005}.  A compiler may target any model-matched forbidden
point from either theorem, but the Williams product region gives the clean
single inequality used here.

Williams's 2024 self-improvement theorem is strong evidence that finite
algorithm-to-algorithm bootstrapping is a legitimate architecture
\cite{Williams2024}.  It is a directly proved result for the circuit-analysis
problems and parameter regimes in that paper.  Its transfer to the present
general-CNF source-localization compiler is only a structural analogy; the
paper does not provide the exact cover, virtual decoder, or charged successor
needed here.

The bounded-treewidth algorithms of Alekhnovich and Razborov and Allender et
al. are directly transferable positive controls for formulas supplied with
the relevant structure \cite{AlekhnovichRazborov2002,AllenderEtAl2014}.  The
missing transfer theorem would have to produce a favorable decomposition or
backdoor on every CNF without paying the parent SAT cost.  Without that
theorem, bounded-width success does not imply a general successor.

The sparsification lemma decomposes fixed-width CNF into a disjunction of
sparse formulas \cite{ImpagliazzoPaturiZane2001}.  It preserves the variable
scale and can use exponentially many disjuncts in the number of variables.
It therefore does not by itself give polynomial fanout $N^a$ of
$N^b$-length children with $b<1$.  Treating sparsification as the missing
wrapper compression is a failed model transfer unless an additional charged
source-size theorem is proved.

\section{Computational exponent audit}

The script \path{pnp_spines/final_link_arc_audit.py} writes
\path{pnp_spines/final_link_arc_audit_results.json}.  It uses exact rational
arithmetic where possible, has an internal 25-second deadline, and completed
14,591 assertions in 3.39 seconds on this run.  The checks included 3,640
small-rational fixed-point cases, four explicit candidate cells, four finite
crossings from starting exponents up to $100$, the PRF3-0079 correction, and
contracting versus unit-product two-type cycles.

For the explicit half-scale example $\beta=1/7$, first entry into $\B$ occurred
after 6, 8, 11, and 13 returns from starting pairs $(2,2)$, $(5,5)$,
$(20,20)$, and $(100,100)$, respectively.  These are arithmetic checks of a
conditional recurrence, not evidence that the compiler exists.

The typed positive control used edges
\[
 (w_{01},\theta_{01})=(1/10,1),\qquad
 (w_{10},\theta_{10})=(1/5,2/5).
\]
The cycle product is $2/5$, and the least time budgets are
$(C_0,C_1)=(1/2,2/5)$.  Replacing the second scale by $1$ produced values
$(3,3)$ after 20 Bellman iterations and $(6,6)$ after 40, illustrating the
unsuppressed unit-product cycle.

\section{Proposed manuscript changes}

No living proof-spine file was edited in this run.  The following changes are
recommended for the next reconciled revision.

\begin{description}
\item[Add a charged successor theorem.]
Insert Theorem 3.1 and criterion \eqref{eq:maincriterion} immediately after
the cofinal-seed theorem.  State that this is a sufficient source-checkable
compiler target, not an established general-SAT construction.

\item[Replace the unique-cell wording.]
Where the spine presents $(a,b)=(1/3,2/3)$ as the target, replace it by the
region
\[
 \max\{r,(a+\beta+g)/(1-b)\}+p<\gamma,
\]
then retain $(1/3,2/3)$ as one example.

\item[Install the PRF3-0079 correction.]
Mark the single-child Route B inequality as superseded for all-block
schedules.  Display $a+\beta+\lambda k$ and, for $a=1-\lambda$, the concrete
constraint $\lambda<1/3-\beta$ at $k=2,C=4/3$.

\item[Add an explicit-versus-virtual representation guard.]
For explicit children require materialization traffic in $r$ and child storage
in $p$.  For virtual children require a bit-level decoder, charge access in
$g$, charge all retained state in $p$, and prove compatibility with the
endpoint machine model.

\item[Add finite portfolio and typed-cycle ledgers.]
Use coordinatewise maxima for portfolios and
Equations~\eqref{eq:belltime}--\eqref{eq:bellspace} for double wrappers.  Allow
same-scale edges only when every reachable return cycle contracts.

\item[Preserve the open status.]
State explicitly that no recovered \textsc{Proved/Exact} component constructs
an all-input strict successor.  The highest-value target is a source-checkable
cover, BatchAcquire object, or virtual decoder whose full charged point lies
strictly inside one verified forbidden region.
\end{description}

\section{Final conclusion}

The final link is smaller than the original fixed wrapper conjecture.  One
does not need to prove that every general-SAT input admits the same
$1/3,2/3$ decomposition.  It is enough to prove either a cofinal family of
finite charged routes, or one exact max-affine compiler whose limiting
time-space envelope lies below the Williams line.  Finite portfolios and
contracting typed return cycles are legitimate alternatives.

The correction is equally important: labels, shared sources, and
source-localized children help only when their global fanout, every access,
and their peak live state are charged.  Under that ledger, the old all-block
half-scale pointwise cell does not survive.  The recovered project work now
supplies a precise acceptance test for the final construction, but it does not
yet supply the construction.  Therefore the present result is a sharpened
conditional proof spine, not a defensible proof that $P\ne NP$.

\begin{thebibliography}{9}

\bibitem{Williams2010}
R. Ryan Williams.
\newblock Alternation-trading proofs, linear programming, and lower bounds.
\newblock In \emph{27th International Symposium on Theoretical Aspects of
Computer Science}, LIPIcs 5, pages 669--680, 2010.
\newblock \url{https://doi.org/10.4230/LIPIcs.STACS.2010.2494}.

\bibitem{FortnowEtAl2005}
Lance Fortnow, Richard J. Lipton, Dieter van Melkebeek, and Anastasios Viglas.
\newblock Time-space lower bounds for satisfiability.
\newblock \emph{Journal of the ACM}, 52(6):835--865, 2005.
\newblock \url{https://doi.org/10.1145/1101821.1101822}.

\bibitem{Williams2024}
R. Ryan Williams.
\newblock Self-improvement for circuit-analysis problems.
\newblock In \emph{Proceedings of the 56th Annual ACM Symposium on Theory of
Computing}, pages 1374--1385, 2024.
\newblock \url{https://doi.org/10.1145/3618260.3649723}.

\bibitem{AlekhnovichRazborov2002}
Michael Alekhnovich and Alexander A. Razborov.
\newblock Satisfiability, branch-width and Tseitin tautologies.
\newblock In \emph{43rd Annual IEEE Symposium on Foundations of Computer
Science}, pages 593--603, 2002.
\newblock \url{https://doi.org/10.1109/SFCS.2002.1181983}.

\bibitem{AllenderEtAl2014}
Eric Allender, Shiteng Chen, Tiancheng Lou, Periklis A. Papakonstantinou, and
Bangsheng Tang.
\newblock Width-parameterized SAT: time-space tradeoffs.
\newblock \emph{Theory of Computing}, 10(12):297--339, 2014.
\newblock \url{https://doi.org/10.4086/toc.2014.v010a012}.

\bibitem{ImpagliazzoPaturiZane2001}
Russell Impagliazzo, Ramamohan Paturi, and Francis Zane.
\newblock Which problems have strongly exponential complexity?
\newblock \emph{Journal of Computer and System Sciences}, 63(4):512--530,
2001.
\newblock \url{https://doi.org/10.1006/jcss.2001.1774}.

\end{thebibliography}

% END CODEX REPORT::LEAD::FINAL-LINK-ARC::2026-10-08
\end{document}
